\exercisesection{1}

\begin{enumerate}
\def\labelenumi{\arabic{enumi}.}
\tightlist
\item
  Let K be a field and A the subring of the ring of formal power series K{[}{[}T{]}{]} consisting of the series in which the coefficient of T is 0; A is a Noetherian local integral domain which is not integrally closed.
\end{enumerate}

\begin{enumerate}
\def\labelenumi{(\alph{enumi})}
\item
  Show that every principal ideal of A, which is non-zero and distinct from A, contains a unique element of the form \(T^{n} + \lambda T^{n+1}\) where \(\lambda \in K\) , \(n \geqslant 2\) ; moreover, this ideal contains all the powers \(T^{m}\) where \(m \geqslant n + 2\) . Deduce that the non-principal ideals of A are the ideals \(AT^{n} + AT^{n+1}\) where \(n \geqslant 2\) .
\item
  Show that every fractional ideal of A is divisorial (prove that every non-principal ideal is the intersection of two principal ideals). Describe the structure of the monoid D(A). The only prime ideals of A are the maximal ideal \(m = AT^{2} + AT^{3}\) and (0).
\end{enumerate}

\begin{enumerate}
\def\labelenumi{\arabic{enumi}.}
\setcounter{enumi}{1}
\tightlist
\item
  \begin{enumerate}
  \def\labelenumii{(\alph{enumii})}
  \tightlist
  \item
    In the factorial polynomial domain K{[}X, Y{]} in two indeterminates over a field K, give an example of two principal ideals a, b such that \(a + b\) is not divisorial.
  \end{enumerate}
\end{enumerate}

\begin{enumerate}
\def\labelenumi{(\alph{enumi})}
\setcounter{enumi}{1}
\item
  Let K be the quadratic extension of the field \(\mathbf{Q}(\mathbf{X})\) of rational functions in one indeterminate over Q, obtained by adjoining to \(\mathbf{Q}(\mathbf{X})\) a root y of the polynomial \(Y^{2}-2X\in\mathbf{Q}(\mathbf{X})[\mathbf{Y}]\) . Let A be the subring of K generated by Z, X and y; A is an integrally closed Noetherian domain. Show that \(p=AX+Ay\) is a prime ideal of height 1 but that \(p^{2}\) is not divisorial (show that \(X^{-1}p^{2}\) is a prime ideal strictly containing p). Moreover \(p.(A:p)\neq A\) .
\item
  Let A be an integrally closed Noetherian local integral domain, whose maximal ideal m is not of height \(\leqslant1\) ; let p be a prime ideal of A of height 1; for every integer i \textgreater{} 0 let \(a_{i} = p + m^{i}\) ; show that \(\text{div}\left(\bigcap_{i} a_{i}\right)\) is distinct from \(\text{sup}(\text{div}(a_{i}))\) .
\end{enumerate}

\begin{enumerate}
\def\labelenumi{\arabic{enumi}.}
\setcounter{enumi}{2}
\tightlist
\item
  Let A be a Noetherian integral domain and p a prime ideal of height 1.
\end{enumerate}

Show that A: \(p \neq A\) (in the field of fractions of A). (If \(c \neq 0\) is an element of p, note that \(p \in \text{Ass}(A/Ac)\) and deduce that Ac: \(p \neq Ac\) by virtue of Chapter IV, §2, Exercise 30). Is the result still valid for a non-Noetherian integral domain (consider a non-discrete valuation ring of height 1)?

\begin{enumerate}
\def\labelenumi{\arabic{enumi}.}
\setcounter{enumi}{3}
\item
  Let A be a completely integrally closed domain and m a maximal ideal of A. Show that, either m is invertible, or m is not divisorial and A: m = A. (Note that, if m is not invertible, necessarily m. \((\mathbf{A}:\mathfrak{m})^{k} = \mathfrak{m}\) for every integer k \textgreater{} 0.)
\item
  Let A be an integral domain, K its field of fractions and U the subgroup of \(K^{*}\) consisting of the invertible elements of A; consider the group \(K^{*}/U = G\) as ordered by the relation derived by passing to the quotient from the relation \(x^{-1}y \in A\) in \(K^{*}\) (Algebra, Chapter VI, § 1, no. 5).
\end{enumerate}

\begin{enumerate}
\def\labelenumi{(\alph{enumi})}
\item
  For a fractional ideal \(\mathfrak{a} \neq 0\) of \(\mathbf{K}\) to be divisorial, it is necessary and sufficient that the image of \(\mathfrak{a} - \{0\}\) in \(\mathbf{G}\) be a major set in \(\mathbf{G}\) (Algebra, Chapter VI, § 1, Exercise 30); for every fractional ideal \(\mathfrak{a} \neq 0\) of \(\mathbf{K}\) , the image of \(\mathfrak{a} - \{0\}\) is the least major set containing the image of \(\mathfrak{a} - \{0\}\) . If \(\mathbf{A}\) is a Krull domain, \(\mathbf{G}\) is an ordered group isomorphic to a directed subgroup of an ordered group of the form \(\mathbf{Z}^{(I)}\) .
\item
  Show that for a homomorphism \(w\) of \(G\) to a totally ordered group \(\Gamma\) to be such that the restriction to \(A - \{0\}\) of the composite mapping \(K^* \to K^*/U \xrightarrow{w} \Gamma\) is a valuation, it is necessary and sufficient that \(w\) be an increasing homomorphism.
\item
  Let H be the directed subgroup of the product ordered group \(Z \times Z\) consisting of the ordered pairs \((s_{1}, s_{2})\) such that \(s_{1} + s_{2} \equiv 0 \pmod{2}\) . Show that there exists no integral domain A such that \(K^{*}/U\) is an ordered group isomorphic to H (deduce from (b) that such a domain would necessarily be a Krull domain, but Proposition 9 of no. 5 would not hold for this ring (cf.~Exercise 22)).
\end{enumerate}

6. Let A be an integral domain.\par

\begin{enumerate}
\def\labelenumi{(\alph{enumi})}
\item
  For a divisorial ideal \(\mathfrak{a}\) of \(\mathbf{A}\) to be such that \(\operatorname{div}(\mathfrak{a})\) is invertible in \(\mathbf{D}(\mathbf{A})\) , it is necessary and sufficient that \(\mathfrak{a} \colon \mathfrak{a} = \mathbf{A}\) (cf.~Exercise 5 and Algebra, Chapter VI, § 1, Exercise 30). In particular, in order that \(\mathfrak{a} \colon \mathfrak{a} = \mathbf{A}\) for every divisorial ideal \(\mathfrak{a}\) of \(\mathbf{A}\) , it is necessary and sufficient that \(\mathbf{A}\) be completely integrally closed.
\item
  In order that \(\mathfrak{a} \colon \mathfrak{a} = \mathbf{A}\) for every finitely generated ideal \(\mathfrak{a} \neq 0\) of \(\mathbf{A}\) , it is necessary and sufficient that \(\mathbf{A}\) be integrally closed.
\end{enumerate}

\begin{enumerate}
\def\labelenumi{\arabic{enumi}.}
\setcounter{enumi}{6}
\item
  Let K be a field and \((v_{\iota})_{\iota \in I}\) a family of discrete valuations on K satisfying conditions \((\mathrm{AK}_{\mathrm{I}})\) and \((\mathrm{AK}_{\mathrm{II}})\) of no. 3 and such that, for every integer r, every family \((n_{h})_{1 \leqslant h \leqslant r}\) of rational integers and every family \((\iota_{h})_{1 \leqslant h \leqslant r}\) of distinct elements of I, there exists \(x \in K\) such that \(v_{\iota_{h}}(x) = n_{h}\) for \(1 \leqslant h \leqslant r\) and \(v_{\iota}(x) \geqslant 0\) for \(\iota\) distinct from the \(\iota_{h}\) . Show that the intersection of the valuation rings of the \(v_{\iota}\) is a Krull domain A whose field of fractions is K and that \((v_{\iota})_{\iota \in I}\) is the family of essential valuations of A. (Show first that K is the field of fractions of A and hence that A is a Krull domain; then prove that for all \(\iota \in I\) the prime ideal of A consisting of those x such that \(v_{\iota}(x) > 0\) is necessarily of height 1, arguing by reductio ad absurdum and using Corollary 2 to Proposition 6 of no. 4.)
\item
  Let A be a Krull domain; show that, for every set I, the polynomial ring \(A[X_{i}]_{i\in I}\) is a Krull domain. (Observe that, if B is a Krull domain, the valuations induced on B by the essential valuations of \(B[X_{1},\ldots,X_{n}]\) are the essential valuations of B.)
\end{enumerate}

¶ 9. (a) Let B be a discrete valuation ring, A = B{[}{[}X{]}{]} the ring of formal power series in one indeterminate over B and C = \(A_{x}\) the ring of fractions \(S^{-1}A\), where S is the multiplicative set of \(X^{h}\) (h ≥ 0). Let v be a normed valuation on B; for every element f = ∑n≥h \(b_n X^{n}\) ≠ 0 of C with \(b_h\) ≠ 0 in B (h positive or negative), we write s(f) = v(\(b_h\) ). Show that s is a Euclidean stathm on C (Algebra, Chapter VII, § 1, Exercise 7), such that s(fg) = s(f) + s(g) for non-zero f, g in C. (If s(f) = p ≥ s(g) = q and h is the least of the degrees of the terms ≠0 in f, show that there exists u ∈ C such that f - ug = \(X^{h+1} f_1\) where \(f_1\) ∈ C and, by arguing by reductio ad absurdum deduce the existence of a process of ``Euclidean division'' in C.) If s(f) = 0, f is invertible in C.

* Show that A is not a Dedekind domain. *

\begin{enumerate}
\def\labelenumi{(\alph{enumi})}
\setcounter{enumi}{1}
\tightlist
\item
  Deduce from (a) that, if B is a Krull domain, A = B{[}{[}X{]}{]} is also a Krull domain. (If K is the field of fractions of B, note that A is the intersection of K{[}{[}X{]}{]} and a family \((\mathrm{C}_{t})\) of principal ideal domains with the same field of fractions as A and that every element of A is invertible in almost all the \(C_{t}\) .)
\end{enumerate}

\begin{enumerate}
\def\labelenumi{\arabic{enumi}.}
\setcounter{enumi}{9}
\tightlist
\item
  Let A be a Krull domain, K its field of fractions, L a field containing K and \((\mathbf{B}_{\alpha})\) a right directed family of Krull domains containing A and contained in L such that the field of fractions \(L_{\alpha}\) of \(B_{\alpha}\) is a finite algebraic extension of K and \(B_{\alpha}\) is the integral closure of A in \(L_{\alpha}\) . For every essential valuation v of A and all \(\alpha\) , let \(e_{\alpha}(v)\) be the sum of the ramification indices of the valuations on \(B_{\alpha}\) which extend v. Show that for the union of the \(B_{\alpha}\) to be a Krull domain, it is necessary and sufficient that for every essential valuation v of A, the set of \(e_{\alpha}(v)\) be bounded.
\end{enumerate}

¶ 11. A divisor \(d\) in an integral domain \(\mathbf{A}\) is called finitely generated if it is of the form \(\operatorname{div}(a)\) , where \(a\) is a finitely generated fractional ideal (which is not necessarily divisorial; cf.~(b)).

\begin{enumerate}
\def\labelenumi{(\alph{enumi})}
\item
  Show that, if A is a Krull domain, every divisor of D(A) is of the form div(a), where \(a = Ax + Ay\) , x, y being elements \(\neq 0\) of the field of fractions of A (use Corollary 2 to Proposition 9 of no. 5).
\item
  Let K be a field, \(A = K[X_{n}]_{n \in N}\) the ring of polynomials over K in a denumerable set of indeterminates, which is a Krull ring (Exercise 8). Let \(A'\) be the subring of A generated by the unit element and the monomials \(X_{i}X_{j}\) for all pairs \((i,j)\) (equivalently, \(A'\) is the set of polynomials all whose terms have even degree). If L and \(L'\) are the fields of fractions of A and \(A'\) , show that \(A' = A \cap L'\) , hence \(A'\) is a Krull ring. Let p be the ideal in \(A'\) generated by the products \(X_{0}X_{i}\) for all \(i \geqslant 0\) ; show that p is a prime ideal of height 1 (hence divisorial) but that it is not finitely generated.
\end{enumerate}

\begin{enumerate}
\def\labelenumi{\arabic{enumi}.}
\setcounter{enumi}{11}
\tightlist
\item
  \begin{enumerate}
  \def\labelenumii{(\alph{enumii})}
  \tightlist
  \item
    Let \(\mathbf{A}\) be an integrally closed domain, \(f(\mathbf{X}) = \sum_{i} a_i \mathbf{X}^i\) and \(g(\mathbf{X}) = \sum_{j} b_j \mathbf{X}^j\) be two polynomials in \(\mathbf{A}[\mathbf{X}]\) ; show that, if \(c \in \mathbf{A}\) is such that all the coefficients of \(f(\mathbf{X}) g(\mathbf{X})\) belong to \(\mathbf{A}c\) , then all the products \(a_i b_j\) belong to \(\mathbf{A}c\) (reduce it to the case where \(\mathbf{A}\) is a valuation ring).
  \end{enumerate}
\end{enumerate}

\begin{enumerate}
\def\labelenumi{(\alph{enumi})}
\setcounter{enumi}{1}
\item
  If A is the ring defined in Exercise 1, give an example of two polynomials \(f, g\) of degree 1 in A{[}X{]} and an element \(c \in \mathbf{A}\) for which the conclusion of (a) does not hold (take \(c = \mathbf{T}^4 + \mathbf{T}^5\) ).
\item
  Under the hypotheses of (a), let a, b, c be the ideals of A generated respectively by the coefficients of f, g and fg: show that \(\text{div}(c) = \text{div}(a) + \text{div}(b)\) . Give an example where \(\tilde{c} \neq \tilde{a}\tilde{b}\) . Show that, if c is principal, c = ab and a and b are invertible; if further A is a local ring, a and b are principal.
\end{enumerate}

\begin{enumerate}
\def\labelenumi{\arabic{enumi}.}
\setcounter{enumi}{12}
\tightlist
\item
  Let A be an integral domain and \((p_{i})_{i\in I}\) a family of non-zero elements such that the principal ideals \(Ap_{i}\) are prime; let S be the multiplicative subset, generated by the \(p_{i}\) .
\end{enumerate}

\begin{enumerate}
\def\labelenumi{(\alph{enumi})}
\item
  Show that \(\mathbf{A} = \mathbf{S}^{-1}\mathbf{A}\cap \left(\bigcap_{i}\mathbf{A}_{\mathbf{A}p_i}\right)\) .
\item
  Suppose that every non-empty family of principal ideals of A has a maximal element. Show that each of the rings \(A_{Ap_{t}}\) is a discrete valuation ring (cf.~Chapter VI, §3, no. 5, Proposition 9). Deduce that, if \(S^{-1}A\) is integrally closed (resp. completely integrally closed), A is integrally closed (resp. completely integrally closed). If \(S^{-1}A\) is a Krull domain, show that A is a Krull domain.
\end{enumerate}

\begin{enumerate}
\def\labelenumi{\arabic{enumi}.}
\setcounter{enumi}{13}
\tightlist
\item
  Let A be a Krull domain and S a saturated multiplicative subset of A not containing 0 (Chapter II, § 2, Exercise 1). Show that, if the canonical homomorphism \(\bar{i}\colon\mathrm{C}(\mathrm{A})\to\mathrm{C}(\mathrm{S}^{-1}\mathrm{A})\) (no. 10, Proposition 17) is bijective, S is gene
\end{enumerate}
% semantic-fragments: legacy-input-seam
\relax{}

rated by a family \((p_{\iota})\) of elements such that each \(A p_{\iota}\) is a prime ideal (note that every divisorial prime ideal which meets S is principal).

¶ 15. (a) Let A be an integral domain and a, b two elements ≠0 of A such that \(Aa \cap Ab = Aab\) . Show that in the polynomial ring A{[}X{]} the principal ideal generated by \(aX + b\) is prime (prove that every polynomial \(f(X) \in A[X]\) such that \(f(-b/a) = 0\) is a multiple of \(aX + b\) , arguing induction on the degree of f).

\begin{enumerate}
\def\labelenumi{(\alph{enumi})}
\setcounter{enumi}{1}
\item
  Let A be a Krull domain and \(a, b\) two elements of A such that \(Aa\) and \(Aa + Ab\) are prime and distinct. Show that \(\mathrm{A}[\mathrm{X}] / (a\mathrm{X} + b)\) is a Krull domain and that \(\mathrm{C}(\mathrm{A}[\mathrm{X}] / (a\mathrm{X} + b))\) is isomorphic to \(\mathrm{C}(\mathrm{A})\) . (Note, with the aid of (a), that \(\mathrm{A}[\mathrm{X}] / (a\mathrm{X} + b)\) is isomorphic to \(\mathrm{A}[b / a] = \mathrm{B}\) ; show that \(Ba\) is prime in B. Observe that \(\mathrm{A}[a^{-1}] = \mathrm{B}[a^{-1}]\) is a Krull domain and use Exercise 12.)
\item
  Let A be an integrally closed Noetherian domain and a, b two elements of the Jacobson radical r of A. Suppose that the ideal \(Aa + Ab\) is a non-divisorial prime ideal; show that Aa and Ab are prime ideals. (Prove first that \(Aa \cap Ab = Aab\) considering the ideals of \(\operatorname{Ass}(A/Aa)\) ). Show secondly that the relations \(xy \in Ab\) and \(y \notin Aa + Ab\) imply \(x \in Ab + Aa^{h}\) for all h, by induction on h; deduce that then \(x \in Ab\) . Finally, deduce that, if \(xy \in Ab\) , \(x \notin Ab\) , \(y \notin Ab\) , then necessarily \(x \in Ab + Aa^{h}\) and \(y \in Ab + Aa^{h}\) for all h, arguing by induction on h; deduce a contradiction.)
\end{enumerate}

\begin{enumerate}
\def\labelenumi{\arabic{enumi}.}
\setcounter{enumi}{15}
\tightlist
\item
  Let \(A = \bigoplus_{n \in Z} A_n\) be a graded Krull domain. Let \(\mathrm{D}_g(\mathrm{A})\) (resp. \(\mathrm{F}_g(\mathrm{A})\) ) denote the group generated by the divisors of the graded integral divisorial ideals (resp. by the divisors \(\operatorname{div}(a)\) , where a is homogeneous in A); write \(\mathrm{C}_g(\mathrm{A}) = \mathrm{D}_g(\mathrm{A}) / \mathrm{F}_g(\mathrm{A})\) . Let S be the multiplicative subset of homogeneous elements \(\neq 0\) of A.
\end{enumerate}

\begin{enumerate}
\def\labelenumi{(\alph{enumi})}
\item
  Show that every prime ideal of A of height 1 which meets S is graded (use Chapter III, § 1, no. 4, Proposition 5).
\item
  Let \(\mathbf{B} = \mathbf{S}^{-1}\mathbf{A}\) , which is a graded ring (Chapter II, § 2, no. 9); write \(\mathbf{B} = \bigoplus_{n\in \mathbf{Z}}\mathbf{B}_n\) ; show that \(\mathbf{B}_0\) is a field and that, if \(\mathbf{A}\neq \mathbf{A}_0\) , \(\mathbf{B}\) is isomorphic to the ring \(\mathbf{B}_0[\mathbf{X},\mathbf{X}^{-1}]\) of rational functions \(\mathrm{P}(\mathbf{X}) / \mathbf{X}^k\) , where \(\mathrm{P}(\mathbf{X})\in \mathrm{B}_0[\mathbf{X}]\) .
\item
  Show that \(C_{g}(A)\) and \(C(A)\) are canonically isomorphic (use (a), (b) and equation (5) of no. 10).
\end{enumerate}

\begin{enumerate}
\def\labelenumi{\arabic{enumi}.}
\setcounter{enumi}{16}
\tightlist
\item
  Let \(\mathbf{A} = \bigoplus_{n\in \mathbf{Z}}\mathbf{A}_n\) be a graded Krull domain and \(p\in \mathbf{A}_1\) be such that \(\mathbf{A}p\) is prime and \(\neq 0\) .
\end{enumerate}

\begin{enumerate}
\def\labelenumi{(\alph{enumi})}
\item
  Show that, if B is the ring \(\mathbf{A}_{(p)}\) defined in Chapter III, § 1, Exercise 1 (a), B is a Krull domain and the groups C(A) and C(B) are canonically isomorphic. (Show first that C(B) is isomorphic to C(B{[}p,p \(^{-1}\) {]}) and observe that B{[}p,p \(^{-1}\) {]} = A{[}p \(^{-1}\) {]}.)
\item
  Under the hypotheses of Exercise 15 (b), show that the groups \(\mathbf{C}(\mathbf{A})\) and \(\mathbf{C}(\mathbf{A}[\mathbf{X},\mathbf{Y}] / (a\mathbf{X} + b\mathbf{Y}))\) are isomorphic.
\end{enumerate}

¶ 18. Let \(A = \bigoplus_{n \in \mathbf{Z}} A_n\) be a graded Krull domain with positive degrees, such that \(A_0\) is a field and let \(m = \bigoplus_{n \geqslant 1} A_n\) , a maximal ideal of \(A\) . Let \(S\) be the multiplicative set of homogeneous elements \(\neq 0\) of \(A\) , so that the ring \(S^{-1}A\) is a principal ideal domain (Exercise 16 (b)).

\begin{enumerate}
\def\labelenumi{(\alph{enumi})}
\item
  Let p be a prime ideal of A of height 1 which meets A --- m; then p is not graded (otherwise p = A) and the ideal S \(^{-1}\) p of S \(^{-1}\) A is principal, generated by an element of the form x = 1 + x \(_{1}\) + · · · + x \(_{n}\) where x \(_{j}\) ∈ (S \(^{-1}\) A) \(_{j}\) . Writing an element of p of the form 1 + a \(_{1}\) + · · · + a \(_{q}\) (a \(_{i}\) ∈ A \(_{i}\) ) as a multiple of x in S \(^{-1}\) A and using Chapter V, § 1, no. 3, Proposition 11, show that x ∈ p; writing finally every element of p as a multiple of x in S \(^{-1}\) A, show that p = Ax.
\item
  Deduce from (a) that C(A) and C(A \(_{m}\) ) are canonically isomorphic (use Proposition 17 of no. 10).
\end{enumerate}

\begin{enumerate}
\def\labelenumi{\arabic{enumi}.}
\setcounter{enumi}{18}
\tightlist
\item
  Let A be a local Krull domain and m its maximal ideal; if \(\mathbf{A}' = (\mathbf{A}[\mathbf{X}])_{\mathfrak{m}\mathbf{A}[\mathbf{X}]}\) , show that \(\mathbf{C}(\mathbf{A})\) and \(\mathbf{C}(\mathbf{A}')\) are canonically isomorphic. (Apply the criterion of Proposition 17 of no. 10, using Exercise 12 (c).)
\end{enumerate}

¶ 20. An integral domain A is called Bezoutian (or a Bezout domain) if every finitely generated ideal of A is principal. Every Bezoutian Noetherian domain is a principal ideal domain. Every valuation ring is a Bezoutian domain (and therefore a Bezoutian domain is not therefore necessarily Noetherian nor completely integrally closed).

\begin{enumerate}
\def\labelenumi{(\alph{enumi})}
\item
  Show that every Bezoutian domain is integrally closed (cf.~Exercise 6). If an integral domain is the union of a right directed family of Bezoutian subdomains, it is Bezoutian. If A is a Bezoutian domain, so is S \(^{-1}\) A for every multiplicative subset S of A such that 0 \(\notin\) S.
\item
  Let \(v_{i}\) ( \(1 \leqslant i \leqslant n\) ) be independent valuations on a field \(K\) and \(A_{i}\) the ring of the valuation \(v_{i}\) . Show that the intersection \(A\) of the \(A_{i}\) is a Bezoutian domain. (If an ideal \(a\) of \(A\) is finitely generated, the set of \(v_{i}(x)\) for \(x \in a\) admits a least element \(\alpha_{i}\) in the value group of \(v_{i}\) ; for all \(i\) , let \(x_{i} \in a\) be such that \(v_{i}(x_{i}) = \alpha_{i}\) . Using the Approximation Theorem (Chapter VI, §7, no. 2,
\end{enumerate}

Theorem 1), show that there are elements \(a_i \in \mathbf{A}\) such that \(x = \sum_{i=1}^{n} a_i x_i \in \mathfrak{a}\) satisfies the relations \(v_i(x) = \alpha_i\) for \(1 \leqslant i \leqslant n\) . If the \(v_i\) are discrete valuations, \(\mathbf{A}\) is a principal ideal domain.

\begin{enumerate}
\def\labelenumi{(\alph{enumi})}
\setcounter{enumi}{2}
\tightlist
\item
  Let K be an algebraically closed field of characteristic \(\neq2\) and let A be the subring of an algebraic closure of \(\mathbf{K}(\mathbf{X})\) generated by \(\mathbf{K}\) and two sequences of elements \((x_{n})\) , \((1 / x_{n})\) ( \(1 \leqslant n < +\infty\) ), where \(x_{1} = \mathbf{X}\) and \(x_{n-1} = x_{n}^{2}\) . Show that \(\mathbf{A}\) is Bezoutian (use (a)). If \(\mathfrak{p}\) is a prime ideal \(\neq 0\) of \(\mathbf{A}\) , show that \(\mathfrak{p}\) is generated by a sequence of elements of the form \(x_{n} - a_{n}\) , where \(a_{n} \in \mathbf{K}\) , \(a_{n} \neq 0\) and \(a_{n}^{2} = a_{n-1}\) (consider for all \(n\) the intersection \(\mathfrak{p} \cap \mathbf{K}[x_{n}, 1 / x_{n}]\) ). Show that \(\mathfrak{p}\) is not finitely generated.
\end{enumerate}

\begin{enumerate}
\def\labelenumi{\arabic{enumi}.}
\setcounter{enumi}{20}
\tightlist
\item
  An integral domain A is called pseudo-Bezoutian (resp. pseudo-principal) if, in the notation of Exercise 5, the group K*/U is lattice-ordered (completely lattice ordered). Every Bezoutian * (resp. factorial) * domain is pseudo-Bezoutian * (resp. pseudo-principal) *. Every pseudo-principal domain is pseudo-Bezoutian. A valuation ring whose order group is R is pseudo-principal but not a principal ideal domain. Every pseudo-Bezoutian (resp. pseudo-principal) domain is integrally closed (resp. completely integrally closed) (use Exercise 6). * Every Noetherian pseudo-Bezoutian domain is factorial. * Give an example of an integrally closed Noetherian (and therefore Krull) domain which is not pseudo-Bezoutian (cf.~Exercise 2). If A is pseudo-Bezoutian, so is S \(^{-1}\) A for every multiplicative subset S of A such that 0∉S (use Chapter II, § 2, Exercise 1).
\end{enumerate}

¶ 22. (a) Let \(\Gamma\) be a lattice-ordered additive group and A the algebra of \(\Gamma\) over a field \(k\) ; A is an integral domain (Algebra, Chapter II, §11, no. 4, Proposition 8). Every element \(x \neq 0\) in A may be written uniquely as \(x = \sum_{i=1}^{n} \alpha_i e^{\nu_i}\) , where the \(\nu_i\) are distinct elements of \(\Gamma\) , the \(e^{\nu_i}\) the corresponding elements of the canonical basis of A over \(k\) and the \(\alpha_i\) elements \(\neq 0\) of \(k\) . Write \(\phi(x) = \inf(\nu_1, \ldots, \nu_n)\) in \(\Gamma\) ; show that \(\phi(x + y) \geqslant \inf(\phi(x), \phi(y))\) if \(x, y\) and \(x + y\) are \(\neq 0\) in A and \(\phi(xy) = \phi(x) + \phi(y)\) if \(xy \neq 0\) in A. (To show the second assertion, establish first the following lemma: given a finite family \((\xi_j)\) of elements of \(\Gamma\) , in order that \(\inf_{i} (\xi_j) = 0\) , it is necessary and sufficient that for all \(\eta > 0\) in \(\Gamma\) there exist an index \(j\) and an element \(\zeta \in \Gamma\) such that \(0 < \zeta \leqslant \eta\) and \(\inf(\xi_j, \zeta) = 0\) . Apply this lemma reducing it to the case where \(\phi(x) = \phi(y) = 0\) .)

\begin{enumerate}
\def\labelenumi{(\alph{enumi})}
\setcounter{enumi}{1}
\tightlist
\item
  Deduce from (a) that, if \(\mathbf{K}\) is the field of fractions of \(\mathbf{A},\phi\) can be extended to a homomorphism from \(\mathbf{K}^*\) to \(\Gamma\) (also denoted by \(\phi\) ) such that
\end{enumerate}

\[
\phi (x + y) \geqslant \inf (\phi (x), \phi (y))
\]

if \(x, y\) and \(x + y\) are \(\neq 0\) in \(\mathbf{K}\) . Deduce that, if \(\mathbf{B}\) is the set of \(x \in \mathbf{K}\) such that \(x = 0\) or \(\phi(x) \geqslant 0\) , \(\mathbf{B}\) is a domain whose field of fractions is \(\mathbf{K}\) and such that, if \(\mathbf{U}\) is the group of invertible elements of \(\mathbf{B}\) , \(\mathbf{K}^*/\mathbf{U}\) is isomorphic to \(\Gamma\) (and in particular \(\mathbf{B}\) is a pseudo-Bezoutian domain). Derive examples of a pseudo-Bezoutian domain which is not completely integrally closed, a completely integrally closed pseudo-Bezoutian domain which is not pseudo-principal (cf.

Algebra, Chapter VI, § 1, Exercise 31) * and a pseudo-principal domain which is not factorial. *

* (c) Derive from (b) an example of an integral domain which is the union of a right directed family of factorial subdomains and is completely integrally closed but not pseudo-Bezoutian. (Let \(\theta\) be an irrational number in {[}0, 1{]} and for every integer \(j\) let \(q_j\) be the greatest integer such that \(q_j / 2^j < \theta\) . For all \(j\) define on the product group \(\mathbf{Q} \times \mathbf{Q}\) a lattice-ordered group structure by taking the set \((\mathrm{G}_j)_+\) of elements \(\geqslant 0\) of this group to be the ordered pairs \((\xi, \eta)\) such that \(\xi \geqslant 0\) and \(0 \leqslant \eta \leqslant (q_j 2^{-j}) \xi\) . *

\begin{enumerate}
\def\labelenumi{\arabic{enumi}.}
\setcounter{enumi}{22}
\tightlist
\item
  \begin{enumerate}
  \def\labelenumii{(\alph{enumii})}
  \tightlist
  \item
    Let A be a pseudo-Bezoutian domain (Exercise 21) and K its field of fractions; for every polynomial \(f \in \mathrm{A}[\mathrm{X}]\) a g.c.d. of the coefficients of \(f\) (determined to within an invertible element of A) is called a content of \(f\) . Let \(f, g\) be two polynomials of \(\mathrm{A}[\mathrm{X}]\) ; show that for \(f\) to divide \(g\) in \(\mathrm{A}[\mathrm{X}]\) , it is necessary and sufficient that \(f\) divide \(g\) in \(\mathrm{K}[\mathrm{X}]\) and that a content of \(f\) divide a content of \(g\) (use Exercise 12) (cf.~Exercise 30 (c)).
  \end{enumerate}
\end{enumerate}

\begin{enumerate}
\def\labelenumi{(\alph{enumi})}
\setcounter{enumi}{1}
\item
  Deduce from (a) that, if A is a pseudo-Bezoutian (resp. pseudo-principal) domain, then so is A{[}X{]}. Moreover, if \((a_{t})\) is a finite (resp. any) family of elements \(\neq0\) of A and d a g.c.d. of this family in A, d is also a g.c.d. of this family in A{[}X{]}.
\item
  Deduce from (b) that, if A is pseudo-Bezoutian (resp. pseudo-principal) then \(A[X_{\lambda}]_{\lambda \in L}\) is pseudo-Bezoutian (resp. pseudo-principal) for every family \((\mathbf{X}_{\lambda})_{\lambda \in L}\) of indeterminates.
\end{enumerate}

\begin{enumerate}
\def\labelenumi{\arabic{enumi}.}
\setcounter{enumi}{23}
\item
  Let K be an algebraically closed field and A = K{[}X, Y{]} the polynomial ring in two indeterminates over K, which is a Krull domain * (and even a factorial domain). * For every ordered pair \((\alpha, \beta) \in \mathrm{K}^{2}\) let \(w_{\alpha,\beta}\) be the discrete valuation on A such that, for every polynomial \(f \neq 0\) , \(w_{\alpha,\beta}(f)\) is the least degree of the monomials \(\neq 0\) in \(f(\mathrm{X} + \alpha, \mathrm{Y} + \beta)\) . Show that A is the intersection of the rings of the valuations \(w_{\alpha,\beta}\) , none of which is an essential valuation of A.
\item
  An integrally closed domain \(\mathbf{A}\) is said to be of finite character if there exists a family \((v_{\iota})_{\iota \in I}\) of valuations on the field of fractions of \(\mathbf{A}\) satisfying properties \((\mathrm{AK}_{\Pi})\) and \((\mathrm{AK}_{\mathrm{III}})\) .
\end{enumerate}

\begin{enumerate}
\def\labelenumi{(\alph{enumi})}
\tightlist
\item
  Show that, if A is integrally closed and of finite character, for every element \(x \neq 0\) in A, there can only be a finite number of extremal elements of the ordered group of principal divisors which are \(\leqslant \operatorname{div}(x)\) . (Let \(J \subset I\) be the finite set of indices such that \(v_{\iota}(x) > 0\) . If \(m\) is the number of elements in J, show that there cannot be \(m + 1\) elements \(y_{j}\) ( \(1 \leqslant j \leqslant m + 1\) ) dividing \(x\) and such that the elements \(\operatorname{div}(y_{j})\) are extremal, observing that, for all \(j\) , \(\operatorname{div}(y_{j})\) must be prime to \(\sum_{k \neq j} \operatorname{div}(y_{k})\) .)
\end{enumerate}

* (b) Show that the ring of integral functions of one complex variable (Chapter V, § 1, Exercise 12) is a pseudo-principal domain (Exercise 21) but not a domain of finite character (use (a)). *

¶ 26. Let A be an integral domain and K its field of fractions. A valuation v on K is called essential for A if the ring of v is a local ring \(A_{p}\) at a prime ideal p of A (the intersection of A and the ideal of v).

\begin{enumerate}
\def\labelenumi{(\alph{enumi})}
\item
  Let \(v\) be an essential valuation for \(A\) of height 1 and let \(p\) be the prime ideal \(A\) consisting of those \(x\) such that \(v(x) > 0\) . Show that \(p\) is of height 1. If \(x \in K\) does not belong to \(A_p\) , then \(A \cap x^{-1}A \subset p\) . If \(w\) is a valuation on \(K\) whose ring \(B\) contains \(A\) and whose ideal \(q\) satisfies \(q \cap A \subset p\) , show that \(w\) is equivalent to \(v\) .
\item
  An integrally closed domain A is said to be of finite character and real type if there exists a family \((v_{\iota})_{\iota\in I}\) of valuations on K of height 1 satisfying properties \((\mathbf{AK}_{\mathrm{II}})\) and \((\mathbf{AK}_{\mathrm{III}})\) . Show that under these conditions, if \(z\in K\) does not belong to A and p is a prime ideal of A such that \(A\cap z^{-1}A\subset p\) , there exists \(\iota\in I\) such that \(v_{\iota}(z)<0\) and that the prime ideal \(q_{\iota}\) of \(x\in A\) such that \(v_{\iota}(x)>0\) is contained in p.~(Argue by reductio ad absurdum by considering the finite number of \(v_{\iota_{k}}\) such that \(v_{\iota_{k}}(z)<0\) ; prove the existence of an \(a\in A\) such that \(a\notin p\) and \(v_{\iota_{k}}(a)>0\) for all k; deduce that \(a^{n}z\notin A\) for every integer n\textgreater0 and show that this implies a contradiction.) Conclude from this that every essential valuation for A of height 1 is equivalent to one of the \(v_{\iota}\) (use (a)). Every finite intersection of subdomains of K which are of finite character and real type is also a domain of finite character and real type.
\item
  Suppose that the hypotheses of (b) are satisfied and moreover that all the \(v_{\iota}\) are essential. Show that, for every prime ideal p of A of height 1, \(A_{p}\) is the ring of one of the valuations \(v_{\iota}\) (use (b) taking \(z^{-1} \in p\) ). For all \(x \in A\) , the principal ideal Ax then admits a unique reduced primary decomposition (Chapter IV, §2, Exercise 20), the prime ideals corresponding to this decomposition being the prime ideals of height 1 containing x.
\item
  Suppose that the hypotheses of (b) are satisfied. Let S be a multiplicative subset of A not containing 0 and \(J \subset I\) the set of \(\iota \in I\) such that \(v_{\iota}(x) = 0\) on S. Show that the family \((v_{\iota})_{\iota \in I-J}\) satisfies properties \((\mathrm{AK}_{\mathrm{II}})\) and \((\mathrm{AK}_{\mathrm{III}})\) for the ring \(S^{-1}A\) ; this family consists of essential valuations for \(S^{-1}A\) if \((v_{\iota})_{\iota \in I}\) consists of essential valuations for A.
\item
  Generalize Proposition 9 of no. 5 to the case where the hypotheses of (c) are satisfied.
\item
  Suppose the hypotheses of (b) are satisfied. Let \(K'\) be a finite extension of K and \(A'\) the integral closure of A in \(K'\) . Show that the valuations (no two of which are equivalent) on \(K'\) which extend the \(v_{\iota}\) satisfy properties ( \(AK_{II}\) ) and ( \(AK_{III}\) ) for \(A'\) and the latter is therefore a domain of finite character and real type; if the \(v_{\iota}\) are all essential for A, so are their extensions for \(A'\) (argue as in no. 8, Proposition 12 and use also Chapter VI, § 8, no. 3, Remark).
\item
  If \(\mathbf{A}\) is a domain of finite character and real type, so is \(\mathbf{A}[\mathbf{X}]\) ; if the conditions of (c) are fulfilled, determine the essential valuations for A{[}X{]} (argue as in no. 9). Generalize Exercise 8 similarly.
\end{enumerate}

\begin{enumerate}
\def\labelenumi{\arabic{enumi}.}
\setcounter{enumi}{26}
\tightlist
\item
  In Exercise 22 take \(\Gamma\) to be a direct sum of a family \((\Gamma_{\iota})_{\iota\in I}\) , where the \(\Gamma_{\iota}\) are subgroups of the additive group \(\mathbf{R}\) . Show that the ring \(\mathbf{B}\) defined in Exercise 22 (b) is a domain of finite character and real type and that it is the intersection of a family of essential valuation rings for \(\mathbf{B}\) ; moreover, every prime ideal \(\neq 0\) of \(\mathbf{B}\) is maximal.
\end{enumerate}

¶ 28. An integrally closed domain A is said to be of finite character and rational type if there exists a family \((v_{\iota})_{\iota \in I}\) of valuations on its field of fractions K satisfying properties \((\mathbf{AK}_{\mathrm{II}})\) and \((\mathbf{AK}_{\mathrm{III}})\) and whose value groups are subgroups of the additive group Q. Show that the family of essential valuations for A of height 1 satisfy \((\mathbf{AK}_{\mathrm{II}})\) and \((\mathbf{AK}_{\mathrm{III}})\) . (For all \(\iota \in I\) , let \(p_{\iota}\) be the prime ideal of A consisting of the x such that \(v_{\iota}(x) > 0\) and let \(V_{\iota}\) be the ring of the valuation \(v_{\iota}\) . Show that, if there are two indices \(\alpha, \beta\) such that \(p_{\beta} \subset p_{\alpha}\) , then A is the intersection of the \(V_{\iota}\) of index \(\iota \neq \alpha\) . For this, argue by reductio ad absurdum by showing that otherwise there would be an element \(x \in K^{*}\) , an element \(y \in p_{\beta}\) and two positive integers r, s such that \(v_{\beta}(x^{r}y^{s}) > 0\) , \(v_{\alpha}(x^{r}y^{s}) = 0\) and \(v_{\iota}(x^{r}y^{s}) \geqslant 0\) for all \(\iota \in I\) , contrary to the hypothesis.)

\begin{enumerate}
\def\labelenumi{\arabic{enumi}.}
\setcounter{enumi}{28}
\tightlist
\item
  Let A be an integrally closed domain and K its field of fractions.
\end{enumerate}

\begin{enumerate}
\def\labelenumi{(\alph{enumi})}
\item
  Let L be an algebraic extension of K and B the integral closure of A in L. Show that for every fractional ideal a of A, if we write b = aB, \(\tilde{b} \cap A = \tilde{a}\) . (Reduce it to the case where \(A \subset \tilde{a}\) , in other words A: \(a \subset A\) and prove that B: \(b \subset B\) ; for this, let \(x \in B\) : b and let \(c_i\) ( \(1 \leqslant i \leqslant n\) ) be the coefficients of its minimal polynomial over K; note that for all \(y \in a\) the elements \(c_i y^i\) ( \(1 \leqslant i \leqslant n\) ) belong to A and deduce that the \(c_i\) belong to A.)
\item
  Let C be a polynomial ring \(\mathbf{A}[\mathbf{X}_{\lambda}]_{\lambda \in L}\) in any family of indeterminates. Show that for every fractional ideal \(\mathfrak{a}\) of A, if we write \(\mathfrak{c} = \mathfrak{aC}\) , \(\tilde{\mathfrak{c}} \cap \mathbf{A} = \tilde{\mathfrak{a}}\) (same method).
\end{enumerate}

¶ 30. An integral domain is called regularly integrally closed if, in the monoid D(A), every finitely generated divisor (Exercise 11) is a regular element. Every completely integrally closed domain is regularly integrally closed; every pseudo-Bezoutial domain (Exercise 21) is regularly integrally closed.

\begin{enumerate}
\def\labelenumi{(\alph{enumi})}
\item
  If A is regularly integrally closed and \(d, d', d''\) are three finitely generated divisors, the relation \(d + d'' \leqslant d' + d''\) implies \(d \leqslant d'\) .
\item
  Deduce from (a) that, for an integral domain A to be regularly integrally closed, it is necessary and sufficient that, for every fractional ideal a of A such that \(\text{div}(a)\) is a finitely generated divisor, \(a: a = A\) ; in particular A is integrally closed (Exercise 6 (b)). (Use the fact that A: (bc) = (A: b): c for two fractional ideals b, c of A, taking b = a, c = A: a.) Moreover, \(\text{div}(a)\) is then invertible in D(A).
\item
  Let A be a regularly integrally closed domain and K its field of fractions; the monoid \(\mathrm{D}_{f}(\mathrm{A})\) of finitely generated divisors of A generates in \(\mathrm{D}(\mathrm{A})\) a lattice-ordered group \(\mathrm{G}_{f}(\mathrm{A})\) . For every polynomial \(p \in K[X]\) , let \(d(p)\) denote the divisor of the fractional ideal of A generated by the coefficients of p; for every rational function r = p/q of \(\mathrm{K}(\mathrm{X})\) , where p, q are in \(K[X]\) and \(q \neq 0\) , \(d(p) - d(q)\) is well defined on \(\mathrm{G}_{f}(\mathrm{A})\) (independently of the expression of r as a quotient of two polynomials); it is denoted by \(d(r)\) (cf.~Exercise 12 (c)). If we then write \(\gamma(r) = ((r), d(r))\) , where (r) is the fractional principal ideal of \(K[X]\) generated by r, \(\gamma\) is an isomorphism of the ordered group \(\mathcal{P}^{*}(\mathrm{A}[\mathrm{X}])\) (notation of §3, no. 2) onto a subgroup of the product ordered group
\end{enumerate}

\[
\mathscr {P} ^ {*} (\mathrm{K} [ \mathrm{X} ]) \times \mathrm{G} _ {f} (\mathrm{A}) \subset \mathscr {P} ^ {*} (\mathrm{K} [ \mathrm{X} ]) \times \mathrm{D} (\mathrm{A}).
\]

Deduce that \(\mathbf{A}[\mathbf{X}]\) is a regularly integrally closed domain and that \(\mathrm{D}_f(\mathbf{A}[\mathbf{X}])\) is isomorphic to \(\mathcal{P}^* (\mathbf{K}[\mathbf{X}])\times \mathrm{D}_f(\mathbf{A})\) (use (a) and (b) and Exercise 29 (b)).

\begin{enumerate}
\def\labelenumi{(\alph{enumi})}
\setcounter{enumi}{3}
\tightlist
\item
  Let B be an integrally closed domain, K its field of fractions and A the subring of the polynomial ring K{[}X, Y{]} consisting of the polynomials whose constant term belongs to B. Show that, if \(B \neq K\) , A is integrally closed but not regularly integrally closed (show that the divisor \(\inf(\text{div}(X), \text{div}(Y))\) corresponds to a divisorial ideal a such that \(a: a \neq A\) ).
\end{enumerate}

¶ 31. (a) Let A be a regularly integrally closed domain (Exercise 30) and K its field of fractions; for every polynomial P ∈ K{[}X{]}, let c(P) denote the divisor div(a), where a is the fractional ideal of K generated by the coefficients of P. Let B denote the subring of K(X) consisting of the rational functions P/Q such that c(P) ≥ c(Q) (show that this condition depends only on the element P/Q, using Exercise 30 (a) and 12 (c)). Then B ∩ K = A.

\begin{enumerate}
\def\labelenumi{(\alph{enumi})}
\setcounter{enumi}{1}
\item
  Show that B is a Bezoutian domain (Exercise 20) (observe that, if P, Q are two polynomials of A{[}X{]}, P + X \(^{m}\) Q divides P and Q in B if m is sufficiently large).
\item
  Show that, if A is a Krull domain, B is a principal ideal domain (consider an increasing sequence of finitely generated fractional principal ideals of B). Give an example where A is completely integrally closed but B is not a principal ideal domain (cf.~Exercise 22).
\item
  Deduce from (c) an example of a Noetherian domain B for which there exists a subfield K of its field of fractions such that \(B \cap K\) is not Noetherian.
\end{enumerate}

\begin{enumerate}
\def\labelenumi{\arabic{enumi}.}
\setcounter{enumi}{31}
\tightlist
\item
  Let A be an integral domain \((a_{i})_{1 \leqslant i \leqslant n}\) a finite family of elements of A, a the ideal \(\sum_{i} A a_{i}\) and R the submodule of \(A^{n}\) generated by the element \((a_{1}, \ldots, a_{n})\) . Show that, for the torsion submodule of the A-module \(M = A^{n}/R\) to be a direct factor of M, it is necessary and sufficient that
\end{enumerate}

\[
\mathfrak {a} (\mathrm{A}: \mathfrak {a}) + (\mathrm{A}: (\mathrm{A}: \mathfrak {a})) = \mathrm{A}.
\]

